% Generated draft: requires mathematical and source review.
\documentclass[11pt]{article}
\usepackage[margin=25mm]{geometry}
\usepackage{fontspec}
\setmainfont{Latin Modern Roman}
\usepackage{amsmath,amssymb,mathtools,unicode-math}
\setmathfont{Latin Modern Math}
\usepackage{enumitem,needspace,xurl,hyperref}
\setlength{\parindent}{0pt}
\setlength{\parskip}{5pt}
\setlength{\emergencystretch}{3em}
\newcommand{\N}{\mathbb{N}}
\newcommand{\Z}{\mathbb{Z}}
\newcommand{\Q}{\mathbb{Q}}
\newcommand{\R}{\mathbb{R}}
\newcommand{\C}{\mathbb{C}}
\newcommand{\F}{\mathbb{F}}
\newcommand{\A}{\mathbb{A}}
\newcommand{\PP}{\mathbb{P}}
\newcommand{\E}{\mathbb{E}}
\newcommand{\eps}{\varepsilon}
\newcommand{\dd}{\,\mathrm{d}}
\begin{document}
\section*{UnsolvedMath ID: 30001400}
\textbf{Source problem code:} catalogue entry
\subsection{Definitions and mathematical statement}
The supplied record asks whether a closure characterization for catalytic majorization remains valid for multiple-copy majorization. It does not explicitly specify the dimensions or whether probability vectors must be finite. The following gives the finite, fixed-dimension interpretation; this domain qualification requires verification against the full statement of Theorem 1.

For an integer \(d\geq1\), define
\[
\Delta_d=\left\{u\in\R^d:u_i\geq0,\ \sum_{i=1}^d u_i=1\right\},\qquad
\Delta_d^{\circ}=\{u\in\Delta_d:u_i>0\text{ for every }i\}.
\]
For \(u,v\in\Delta_d\), write \(u\prec v\) if their coordinates, arranged in nonincreasing order, satisfy
\[
\sum_{i=1}^r u_i^{\downarrow}\leq\sum_{i=1}^r v_i^{\downarrow}
\quad(1\leq r<d),\qquad
\sum_{i=1}^d u_i=\sum_{i=1}^d v_i.
\]
For finite vectors \(u\) and \(z\), their tensor product has coordinates \((u\otimes z)_{(i,j)}=u_i z_j\). Define
\[
u\prec_C v
\quad\Longleftrightarrow\quad
\exists m\geq1\ \exists z\in\Delta_m:\ u\otimes z\prec v\otimes z,
\]
and
\[
u\prec_M v
\quad\Longleftrightarrow\quad
\exists n\in\Z,\ n\geq1:\ u^{\otimes n}\prec v^{\otimes n}.
\]
Here the tensor power has coordinates \(u_{i_1}\cdots u_{i_n}\). For \(u\in\Delta_d^{\circ}\) and real \(p\), put
\[
N_p(u)=\sum_{i=1}^d u_i^p.
\]
In particular, \(N_0(u)=d\) and \(N_1(u)=1\).

Fix \(x,y\in\Delta_d^{\circ}\). The supplied characterization for \(\prec_C\) asserts that the conditions
\[
\begin{aligned}
N_p(x)&\leq N_p(y)&&\text{for every real }p\geq1,\\
N_p(x)&\geq N_p(y)&&\text{for every real }0\leq p\leq1,\\
N_p(x)&\leq N_p(y)&&\text{for every real }p\leq0
\end{aligned}
\tag{P}
\]
characterize the existence of probability vectors \(x_k\prec_C y\) converging to \(x\) in \(\ell^1\).

The question is whether, under the finite, fixed-dimension interpretation above, for every \(d\geq1\) and every \(x,y\in\Delta_d^{\circ}\), condition \(\mathrm{(P)}\) is equivalent to
\[
\exists (x_k)_{k\geq1}\subseteq\Delta_d\ \exists(n_k)_{k\geq1}\subseteq\Z_{\geq1}:\quad
x_k^{\otimes n_k}\prec y^{\otimes n_k}\ \text{for every }k,
\qquad
\lim_{k\to\infty}\sum_{i=1}^d |(x_k)_i-x_i|=0.
\]
The integer \(n_k\) may depend on \(k\); no common tensor exponent is required. In this fixed-dimension interpretation, convergence to strictly positive \(x\) makes all sufficiently late \(x_k\) strictly positive, so requiring positive approximants gives the same limit condition. If varying dimensions or infinite probability vectors were intended, the ambient space, majorization convention, and treatment of zero coordinates and nonpositive powers would require additional specification. The stated catalytic characterization is supplied source content, not a theorem independently verified in this draft.
\subsection{Short English statement}
For strictly positive probability vectors, the supplied problem asks whether the same three families of power-sum inequalities that characterize limits under catalytic majorization also characterize limits under multiple-copy majorization. Each approximating vector may use its own number of tensor copies. The finite, fixed-dimension interpretation is explicit here but needs source verification. The record labels the problem open; this draft does not independently verify that status.
\subsection{Sources}
\begin{itemize}
\item Local UnsolvedMath registry: supplied record for permanent ID \(30001400\), exact source problem code \texttt{catalogue entry}, title ``Comparison of Majorization Relations for Probability Vectors,'' in the Oberwolfach Reports Open Problems collection. The record gives proposed year \(2009\) and supplies no individual proposer attribution.
\item Supplied report citation: \emph{Mini-Workshop: Geometry of Quantum Entanglement (2010)}, Oberwolfach Reports, pp.\ 2993--3032, DOI \texttt{10.4171/owr/2009/56}. Supplied source URL and literature-source link, labelled ``original Oberwolfach problem report'': \url{https://doi.org/10.4171/owr/2009/56}. The report is described as documenting workshop discussions and open problems. The full report and Theorem 1 were not supplied for direct inspection.
\item Literature assessment checked at \texttt{2026-08-21}: this is supplied dated curation, not independently verified. It labels the problem open and states that the report explicitly leaves open the extension of Theorem 1 to the multiple-copy order and that no later proof was located. The supplied background requests independent verification of time-sensitive status and cited claims.
\item The registry labels the statement \texttt{corrected\_verified} and attributes verification of the definitions, power-sum conditions, and question to the official OWR report. This is supplied verification metadata, not independent verification performed for this draft.
\end{itemize}
\end{document}
